% !TeX root = ../../ONMIPA Alin.tex
% Sumber utama: Zhang (2009), Bab 3; Zhang (2011), bagian 1.3, 2.4, 3.3.
\chapter{Nilai Eigen dan Vektor Eigen}
\label{ch:eigen}

Bab ini menghubungkan akar polinom karakteristik dengan ruang eigen,
rank, dan nulitas. Pembahasan berlanjut ke keserupaan, polinom minimal,
Cayley--Hamilton, serta persamaan matriks. Diagonalisasi dan bentuk-bentuk
dekomposisi dikembangkan tersendiri pada \autoref{ch:dekomposisi}.

Matriks berentri dalam $\F=\R$ atau $\C$; medan yang digunakan disebutkan
secara eksplisit ketika seluruh akar suatu polinom diperlukan.

\section{Nilai Eigen dan Ruang Eigen}
\label{eigen:dasar}

\begin{definisi}[][Nilai eigen dan vektor eigen][eigen:definisi]
    Untuk $A\in\Mat_n(\F)$, skalar $\lambda\in\F$ disebut
    \emph{nilai eigen} jika terdapat $0\ne v\in\F^n$ dengan
    $Av=\lambda v$. Vektor $v$ disebut \emph{vektor eigen} yang
    bersesuaian dengan $\lambda$. Ruang eigennya adalah
    \[
        E_\lambda(A)=\ker(A-\lambda I).
    \]
    Polinom karakteristik $A$ didefinisikan oleh
    $\chi_A(t)=\det(tI-A)$.
\end{definisi}

Ruang eigen memuat vektor nol, tetapi vektor eigen selalu tak nol.
Dengan hasil pada bab rank dan determinan,
\[
    \lambda\text{ nilai eigen}
    \iff\chi_A(\lambda)=0
    \iff\rank(A-\lambda I)<n.
\]
Menghitung nilai eigen berarti menentukan akar polinom karakteristik;
menghitung vektor eigen berarti menyelesaikan sistem homogen untuk
setiap akar tersebut.

\begin{definisi}[][Multiplisitas aljabar dan geometri][eigen:multiplicitas]
    Multiplisitas aljabar $a_\lambda$ adalah banyaknya faktor
    $t-\lambda$ dalam $\chi_A(t)$. Multiplisitas geometri adalah
    \[
        g_\lambda=\dim E_\lambda(A)=\nul(A-\lambda I).
    \]
\end{definisi}

\begin{teorema}[][Batas dimensi ruang eigen][eigen:batas-geometri]
    Untuk setiap nilai eigen $\lambda$, berlaku
    $1\le g_\lambda\le a_\lambda$.
    Ruang-ruang eigen untuk nilai eigen yang berbeda berjumlah langsung.
\end{teorema}
\begin{proof}
    Pilih basis $v_1,\ldots,v_g$ untuk $E_\lambda(A)$ dan perluas
    menjadi basis $\F^n$. Dalam basis ini, matriks $A$ berbentuk
    \[
        \begin{pmatrix}\lambda I_g&C\\0&D\end{pmatrix}.
    \]
    Determinan segitiga blok memberi
    $\chi_A(t)=(t-\lambda)^g\chi_D(t)$, sehingga $g\le a_\lambda$.

    Untuk sifat jumlah langsung, ambil nilai eigen berbeda
    $\lambda_1,\ldots,\lambda_s$ dan relasi
    $x_1+\cdots+x_s=0$ dengan $x_i\in E_{\lambda_i}(A)$.
    Untuk suatu $i$, terapkan $\prod_{j\ne i}(A-\lambda_jI)$.
    Semua suku kecuali $x_i$ lenyap, sehingga
    \[
        \left(\prod_{j\ne i}(\lambda_i-\lambda_j)\right)x_i=0.
    \]
    Koefisiennya tak nol; maka setiap $x_i=0$.
\end{proof}

Khususnya, vektor-vektor eigen untuk nilai eigen berbeda bebas linear.
Namun, beberapa vektor eigen untuk \emph{nilai eigen yang sama}
masih perlu diperiksa kebebasannya.

\begin{prop}[Trace, determinan, dan rank]
    \label{eigen:trace-det}
    Jika nilai eigen $A\in\Mat_n(\C)$ adalah
    $\lambda_1,\ldots,\lambda_n$, dihitung dengan multiplisitas aljabar,
    maka
    \[
        \chi_A(t)=\prod_{i=1}^n(t-\lambda_i),\qquad
        \tr A=\sum_{i=1}^n\lambda_i,\qquad
        \det A=\prod_{i=1}^n\lambda_i.
    \]
    Banyaknya nilai eigen tak nol, dengan multiplisitas, tidak melebihi
    $\rank A$.
\end{prop}
\begin{proof}
    Faktorkan polinom karakteristik di atas $\C$.
    Koefisien $t^{n-1}$ dalam $\det(tI-A)$ adalah
    $-\sum_i a_{ii}=-\tr A$, sedangkan suku konstannya $(-1)^n\det A$.
    Membandingkan koefisien memberi dua identitas.
    Jika $0$ adalah nilai eigen dengan multiplisitas aljabar $a_0$,
    maka $\rank A=n-g_0\ge n-a_0$.
    Jika $0$ bukan nilai eigen, $A$ invertibel dan ranknya $n$.
\end{proof}

Rank dapat lebih besar daripada banyaknya nilai eigen tak nol.
Misalnya, $\begin{pmatrix}0&1\\0&0\end{pmatrix}$ mempunyai rank satu,
tetapi kedua nilai eigennya nol. Ini membedakan informasi kernel dari
informasi akar polinom karakteristik.


\begin{exam}[Nilai eigen matriks rank satu]
    Misalkan $A=uv^T$, dengan $u,v\in\F^n$ tak nol dan $n\ge2$.
    Tuliskan $c=v^Tu$. Tentukan polinom karakteristik, ruang eigen,
    dan pangkat $A$.
\end{exam}
\begin{solution}
    Karena $Ax=u(v^Tx)$, diperoleh $\rank A=1$ dan
    $\ker A=\ker v^T$ berdimensi $n-1$.
    Trace $A$ adalah $c$, sehingga
    \[
        \chi_A(t)=t^{n-1}(t-c).
    \]
    Jika $c\ne0$, ruang eigen untuk $c$ adalah $\operatorname{span}\{u\}$,
    sedangkan ruang eigen untuk $0$ adalah $\ker v^T$.
    Jika $c=0$, satu-satunya nilai eigen adalah $0$, dan ruang eigennya
    tetap berdimensi $n-1$. Vektor $u$ kini juga berada pada kernel.
    Identitas $A^2=u(v^Tu)v^T=cA$ memberi
    $A^k=c^{k-1}A$ untuk $k\ge1$, dengan faktor berpangkat nol diartikan $1$.
\end{solution}

\section{Keserupaan dan Perubahan Basis}
\label{eigen:similaritas}

\begin{definisi}[][Keserupaan matriks][eigen:similar]
    Matriks $A,B\in\Mat_n(\F)$ disebut serupa jika
    $B=P^{-1}AP$ untuk suatu matriks invertibel $P$.
\end{definisi}

Kolom-kolom $P$ menyatakan basis baru dalam koordinat lama.
Persamaan $AP=PB$ berarti kolom ke-$j$ dari $B$ menyimpan koordinat
bayangan vektor basis ke-$j$ terhadap basis baru tersebut.
Keserupaan mempertahankan polinom karakteristik, trace, determinan,
dan rank, sebab
\[
    tI-B=P^{-1}(tI-A)P.
\]
Kesamaan nilai eigen beserta multiplisitasnya belum cukup untuk
menjamin keserupaan.

\section{Polinom Matriks dan Cayley--Hamilton}
\label{eigen:polinom}

Untuk $f(t)=a_0+a_1t+\cdots+a_dt^d$, definisikan
$f(A)=a_0I+a_1A+\cdots+a_dA^d$.
Semua polinom dalam satu matriks $A$ saling komutatif.
Jika $Av=\lambda v$, maka $f(A)v=f(\lambda)v$.

\begin{prop}[Nilai eigen polinom matriks]
    \label{eigen:spectral-mapping}
    Untuk $A\in\Mat_n(\C)$ dan polinom $f$, nilai eigen $f(A)$
    tepat merupakan bilangan $f(\lambda)$ dengan $\lambda$ nilai eigen $A$.
    Khususnya, $f(A)$ invertibel jika dan hanya jika
    $f(\lambda)\ne0$ untuk setiap nilai eigen $A$.
\end{prop}
\begin{proof}
    Satu arah mengikuti $f(A)v=f(\lambda)v$.
    Jika $f$ tidak konstan, faktorkan
    $f(t)-\mu=c\prod_j(t-\alpha_j)$ di atas $\C$.
    Matriks $f(A)-\mu I=c\prod_j(A-\alpha_jI)$ singular tepat ketika
    salah satu faktornya singular. Jadi $\mu=f(\alpha_j)$ untuk
    suatu nilai eigen $\alpha_j$ dari $A$.
    Untuk $f$ konstan, pernyataannya langsung.
\end{proof}

\begin{teorema}[][Cayley--Hamilton][eigen:ch]
    Setiap matriks persegi memenuhi polinom karakteristiknya:
    jika $\chi_A(t)=t^n+c_{n-1}t^{n-1}+\cdots+c_0$, maka
    \[
        A^n+c_{n-1}A^{n-1}+\cdots+c_1A+c_0I=0.
    \]
\end{teorema}
\begin{proof}
    Ekspansikan adjugat sebagai polinom matriks
    $\operatorname{adj}(tI-A)=\sum_{j=0}^{n-1}B_jt^j$.
    Identitas adjugat memberi
    \[
        (tI-A)\sum_{j=0}^{n-1}B_jt^j=\chi_A(t)I.
    \]
    Perbandingan koefisien menghasilkan
    \[
        B_{n-1}=I,\qquad B_{j-1}-AB_j=c_jI\ (1\le j<n),
        \qquad -AB_0=c_0I.
    \]
    Kalikan persamaan untuk $c_j$ dari kiri dengan $A^j$, kemudian
    jumlahkan bersama persamaan konstanta dan $A^n=A^nB_{n-1}$.
    Suku-suku yang memuat $B_j$ saling menghapus, menyisakan
    $\chi_A(A)=0$.
\end{proof}

Identitas ini diperoleh melalui koefisien polinom. Substitusi formal
$t=A$ ke dalam simbol determinan $\det(tI-A)$ bukan pembuktian:
semula entri matriks tersebut adalah polinom skalar, sedangkan
substitusi itu mengganti jenis entrinya dengan matriks.

\begin{exam}[Pangkat melalui rekurensi]
    Untuk $F=\begin{pmatrix}1&1\\1&0\end{pmatrix}$, tentukan $F^k$.
\end{exam}
\begin{solution}
    Karena $\chi_F(t)=t^2-t-1$, Cayley--Hamilton memberi $F^2=F+I$.
    Dengan bilangan Fibonacci $f_0=0$, $f_1=1$, dan
    $f_{k+1}=f_k+f_{k-1}$, induksi memberikan
    \[
        F^k=f_kF+f_{k-1}I
        =\begin{pmatrix}f_{k+1}&f_k\\f_k&f_{k-1}\end{pmatrix}
        \qquad(k\ge1).
    \]
    Jadi setiap pangkat tereduksi menjadi kombinasi $I$ dan $F$.
    Selain itu $F(F-I)=I$, sehingga $F^{-1}=F-I$.
\end{solution}

\subsection{Polinom minimal dan akar berulang}

\begin{definisi}[][Polinom minimal][eigen:minimal]
    Polinom minimal $m_A(t)$ adalah polinom monik berderajat terkecil
    yang memenuhi $m_A(A)=0$.
\end{definisi}

\begin{prop}[Sifat polinom minimal]
    \label{eigen:sifat-minimal}
    Polinom minimal ada dan unik. Jika $f(A)=0$, maka $m_A\mid f$;
    khususnya $m_A\mid\chi_A$.
    Akar-akar $m_A$ tepat nilai eigen $A$, dan matriks serupa mempunyai
    polinom minimal yang sama.
\end{prop}
\begin{proof}
    Keberadaan mengikuti Cayley--Hamilton.
    Bagi $f$ oleh $m_A$: $f=qm_A+r$ dengan $\deg r<\deg m_A$.
    Maka $r(A)=0$, sehingga minimalitas memaksa $r=0$.
    Ini sekaligus membuktikan keunikan polinom monik minimal.

    Jika $Av=\lambda v$ dengan $v\ne0$, maka
    $0=m_A(A)v=m_A(\lambda)v$, sehingga $m_A(\lambda)=0$.
    Sebaliknya, jika $m_A(t)=(t-\lambda)q(t)$ tetapi $A-\lambda I$
    invertibel, $m_A(A)=0$ memaksa $q(A)=0$, bertentangan dengan minimalitas.
    Terakhir, $f(P^{-1}AP)=P^{-1}f(A)P$ menunjukkan bahwa matriks
    serupa mempunyai polinom penganulir yang sama.
\end{proof}


Kriteria polinom minimal untuk diagonalisasi dan penggunaan proyektor
akan dibahas pada \autoref{ch:dekomposisi}.

\section{Hasil Kali dan Persamaan Matriks}
\label{eigen:hasilkali}

\begin{teorema}[][Nilai eigen $AB$ dan $BA$][eigen:abba]
    Untuk $A\in\Mat_{m\times n}(\C)$ dan $B\in\Mat_{n\times m}(\C)$,
    \[
        t^n\chi_{AB}(t)=t^m\chi_{BA}(t).
    \]
    Jadi $AB$ dan $BA$ mempunyai nilai eigen tak nol yang sama,
    termasuk multiplisitas aljabarnya.
\end{teorema}
\begin{proof}
    Untuk $t\ne0$, hitung determinan
    $\begin{pmatrix}tI_m&A\\B&I_n\end{pmatrix}$ dengan dua komplemen Schur.
    Diperoleh
    \[
        \det(tI_m-AB)=t^m\det(I_n-t^{-1}BA).
    \]
    Kalikan dengan $t^n$ untuk mendapatkan identitas yang dinyatakan.
    Kedua ruas adalah polinom yang sama pada semua $t\ne0$,
    sehingga identitas berlaku juga pada $t=0$.
\end{proof}

Jika $ABx=\lambda x$ dengan $\lambda\ne0$, maka $Bx\ne0$ dan
$BA(Bx)=\lambda Bx$. Bahkan $B$ memberi isomorfisme
$E_\lambda(AB)\to E_\lambda(BA)$ dengan invers $A/\lambda$.
Kesamaan multiplisitas aljabar pada teorema memerlukan identitas polinom,
sedangkan argumen ruang eigen ini menjelaskan multiplisitas geometrinya.

Untuk $m=n$, polinom karakteristiknya sama, tetapi matriksnya belum
tentu serupa. Ambil
\[
    A=\begin{pmatrix}1&0\\0&0\end{pmatrix},\qquad
    B=\begin{pmatrix}0&1\\0&0\end{pmatrix}.
\]
Maka $AB=B\ne0$, sedangkan $BA=0$; rank keduanya berbeda.
Jika salah satu faktor invertibel, keserupaan berlaku, misalnya
$BA=A^{-1}(AB)A$ ketika $A$ invertibel.

\begin{prop}[Persamaan Sylvester dengan spektrum terpisah]
    \label{eigen:sylvester}
    Jika $A\in\Mat_m(\C)$ dan $B\in\Mat_n(\C)$ tidak mempunyai
    nilai eigen bersama, maka untuk setiap $C\in\Mat_{m\times n}(\C)$,
    persamaan $AX-XB=C$ mempunyai tepat satu solusi.
\end{prop}
\begin{proof}
    Jika $AX=XB$, maka $A^kX=XB^k$, sehingga
    $\chi_A(A)X=X\chi_A(B)$.
    Ruas kiri nol oleh Cayley--Hamilton; matriks $\chi_A(B)$
    invertibel oleh \autoref{eigen:spectral-mapping} dan pemisahan spektrum.
    Jadi $X=0$.
    Pemetaan linear $X\mapsto AX-XB$ pada ruang berdimensi $mn$
    dengan demikian injektif, sehingga bijektif.
\end{proof}

\section{Latihan Pilihan}
\label{eigen:latihan}

Latihan berikut menghubungkan ruang eigen, polinom, dan persamaan matriks.

\begin{latihan}\label{eigen:lat-polinom}
    Misalkan $A\in\Mat_m(\C)$ dan $B\in\Mat_n(\C)$.
    Buktikan bahwa $\chi_B(A)$ invertibel jika dan hanya jika
    $A$ dan $B$ tidak mempunyai nilai eigen bersama.
\end{latihan}

\begin{latihan}\label{eigen:lat-dua-rankone}
    Misalkan $a=(a_1,\ldots,a_n)^T\in\R^n$, $n\ge2$,
    $\sum_i a_i=0$, dan $e=(1,\ldots,1)^T$.
    Tentukan semua nilai eigen dan ruang eigennya untuk $M=aa^T+ee^T$.
    Tangani juga kasus $a=0$ dan $\|a\|^2=n$.
\end{latihan}

\begin{latihan}\label{eigen:lat-persamaan}
    Misalkan
    \[
        A=\begin{pmatrix}1&1\\0&2\end{pmatrix},\qquad
        B=\begin{pmatrix}2&0\\0&3\end{pmatrix},\qquad
        C=(c_{ij})\in\Mat_2(\C).
    \]
    Tentukan syarat konsistensi dan seluruh solusi $AX-XB=C$.
    Hubungkan parameter bebasnya dengan nilai eigen bersama $A$ dan $B$.
\end{latihan}

\begin{latihan}\label{eigen:lat-jumlah-vektor}
    Misalkan $u,v$ vektor eigen $A$ untuk nilai eigen yang berbeda.
    Buktikan bahwa $u+v$ bukan vektor eigen $A$.
\end{latihan}

\begin{latihan}\label{eigen:lat-trace-komutator}
    Untuk $A,B\in\Mat_n(\C)$, buktikan $\tr(AB)=\tr(BA)$.
    Simpulkan bahwa $AB-BA=I$ tidak mungkin terjadi.
    Selanjutnya, jika $AB-BA=A$, buktikan $A$ singular.
\end{latihan}

\begin{latihan}\label{eigen:lat-adjugat}
    Misalkan $A\in\Mat_n(\C)$ singular, $n\ge2$, dengan
    $\chi_A(t)=t^n+c_{n-1}t^{n-1}+\cdots+c_1t$.
    Buktikan
    \[
        \chi_{\operatorname{adj}A}(t)
        =t^{n-1}\bigl(t-(-1)^{n-1}c_1\bigr).
    \]
\end{latihan}

\ifpembahasan
\section{Pembahasan Latihan}
\label{eigen:pembahasan}

\subsection*{Latihan \ref{eigen:lat-polinom}: akar karakteristik sebagai penghalang invers}
Dengan \autoref{eigen:spectral-mapping}, $\chi_B(A)$ invertibel
tepat ketika $\chi_B(\lambda)\ne0$ untuk semua nilai eigen $\lambda$ dari $A$.
Nol-nol $\chi_B$ tepat merupakan nilai eigen $B$.
Jadi syarat tersebut ekuivalen dengan tidak adanya nilai eigen bersama.
Ukuran $A$ dan $B$ tidak harus sama karena $\chi_B$ dievaluasi sebagai
polinom skalar pada matriks $A$.

\subsection*{Latihan \ref{eigen:lat-dua-rankone}: dua arah yang ortogonal}
Karena $a^Te=0$, diperoleh $Me=ne$ dan $Ma=\|a\|^2a$.
Untuk $x$ ortogonal terhadap $a$ dan $e$, berlaku $Mx=0$.
Jika $a\ne0$ dan $\|a\|^2\ne n$, ruang eigennya adalah
\[
    E_n=\operatorname{span}\{e\},\quad
    E_{\|a\|^2}=\operatorname{span}\{a\},\quad
    E_0=\operatorname{span}\{a,e\}^{\perp}.
\]
Dimensinya $1,1,n-2$; bila $n=2$, ruang terakhir nol dan $0$ bukan nilai eigen.
Jika $a\ne0$ dan $\|a\|^2=n$, kedua ruang pertama bergabung menjadi
$E_n=\operatorname{span}\{a,e\}$, berdimensi dua.
Jika $a=0$, $M=ee^T$ mempunyai nilai eigen $n$ pada
$\operatorname{span}\{e\}$ dan $0$ pada $e^\perp$.

\subsection*{Latihan \ref{eigen:lat-persamaan}: nilai eigen bersama menyisakan parameter}
Tuliskan $X=\begin{pmatrix}x&y\\z&w\end{pmatrix}$.
Perhitungan memberikan
\[
    AX-XB=\begin{pmatrix}-x+z&-2y+w\\0&-w\end{pmatrix}.
\]
Jadi konsistensi tepat ketika $c_{21}=0$. Jika konsisten, semua solusi adalah
\[
    X=\begin{pmatrix}
        t-c_{11}&-(c_{12}+c_{22})/2\\t&-c_{22}
    \end{pmatrix},\qquad t\in\C.
\]
Kedua matriks mempunyai satu nilai eigen bersama, yaitu $2$, dengan
ruang eigen masing-masing berdimensi satu. Satu pasangan arah ini
memberikan kernel berdimensi satu pada pemetaan $X\mapsto AX-XB$.

\subsection*{Latihan \ref{eigen:lat-jumlah-vektor}: kebebasan dua vektor eigen}
Tuliskan $Au=\lambda u$ dan $Av=\mu v$, dengan $\lambda\ne\mu$.
Kedua vektor bebas linear. Jika $A(u+v)=\nu(u+v)$, maka
$(\lambda-\nu)u+(\mu-\nu)v=0$, sehingga
$\lambda=\nu=\mu$, suatu kontradiksi.

\subsection*{Latihan \ref{eigen:lat-trace-komutator}: trace menghalangi invertibilitas}
Ekspansi entri memberi
$\tr(AB)=\sum_{i,j}a_{ij}b_{ji}=\tr(BA)$.
Jadi trace suatu komutator $AB-BA$ selalu nol,
sedangkan $\tr I=n\ne0$.
Sekarang anggap $AB-BA=A$ dan $A$ invertibel.
Mengalikan dari kiri dengan $A^{-1}$ memberi
$B-A^{-1}BA=I$. Trace ruas kiri nol karena
$\tr(A^{-1}BA)=\tr B$, tetapi trace ruas kanan $n$.
Kontradiksi ini membuktikan $A$ singular.

\subsection*{Latihan \ref{eigen:lat-adjugat}: satu nilai eigen yang mungkin tak nol}
Jika $\rank A\le n-2$, semua minor berukuran $n-1$ nol,
sehingga $\operatorname{adj}A=0$.
Jika $\rank A=n-1$, identitas $A\operatorname{adj}A=0$
menempatkan semua kolom adjugat pada kernel $A$ yang berdimensi satu.
Jadi dalam kedua kasus $\rank(\operatorname{adj}A)\le1$.
Matriks berank paling besar satu mempunyai polinom karakteristik
$t^{n-1}(t-\tr M)$: multiplisitas aljabar nol minimal $n-1$,
dan trace menentukan akar yang tersisa.
Koefisien $t$ dalam $\det(tI-A)$, melalui ekspansi minor utama,
adalah $(-1)^{n-1}$ kali jumlah minor utama berukuran $n-1$.
Jumlah itu tepat $\tr(\operatorname{adj}A)$.
Dengan demikian $\tr(\operatorname{adj}A)=(-1)^{n-1}c_1$,
yang menghasilkan rumus yang diminta.
\fi
